% !TEX root = ../paper.tex

\section{Additional Permit Evidence}
\label{sec:appendix_permits}

This appendix reports additional permit evidence that complements the descriptive discussion in Section~\ref{sec:data} and the event study in Section~\ref{sec:empirical_results_permits}.

\subsection{Balance Before Redistricting}

Table~\ref{tab:permit_itt_balance} compares blocks that changed wards in 2015 with nearby blocks that remained in the same ward.
Before the new map took effect, the two groups had similar permit activity and neighborhood characteristics.

\input{../tasks/permit_itt_balance/output/permit_itt_balance_500ft.tex}

\subsection{Permit Event-Study Estimates}

Table~\ref{tab:permit_event_study_did} reports pooled estimates for years 0 through 5 after redistricting from the signed model used in the main text and from a model that estimates the two directions separately.
Low-discretion permits show a more muted response.

\begin{table}[htbp]
    \centering
    \caption{Signed and Unconstrained Reassignment Estimates}
    \label{tab:permit_event_study_did}
    \input{../tasks/audits/permit_event_study_audit/output/binary_event_study_appendix_income_added_back_500ft.tex}

    \footnotesize
    \textit{Notes:} PPML estimates use census-block-year observations from 2010--2020 within 500ft of a ward boundary.
    Stringency scores are estimated using permits filed from 2006 through 2014.
    Blocks are classified as assigned toward greater stringency, assigned toward greater leniency, or unchanged under the new map.
    The signed-direction effect constrains assignments toward greater stringency and greater leniency to have equal-and-opposite effects in logs.
    The unconstrained model estimates the two assignment effects separately.
    The directional contrast equals one-half of the difference between the more-stringent and more-lenient coefficients.
    The symmetry test evaluates whether the two unconstrained coefficients sum to zero.
    Each model includes block fixed effects, ward-pair $\times$ year fixed effects, and the outcome's pre-period volume and zero-volume indicator interacted with year.
    Fixed-effect PPML drops blocks with no outcome in any sample year and ward-pair-year groups with no outcome.
    Standard errors, in parentheses, are clustered by ward pair.
    $^{*}p<0.10$, $^{**}p<0.05$, $^{***}p<0.01$.
\end{table}

\subsection{Unconstrained Assignment Effects}
\label{sec:appendix_permit_unconstrained}

Figure~\ref{fig:permit_event_study_unconstrained} relaxes the symmetry restriction imposed by the signed model.
The resulting directional contrast is nearly identical to the main estimate.

\begin{figure}[htbp]
    \centering
    \includegraphics[width=0.7\textwidth]{../tasks/audits/permit_event_study_audit/output/binary_joint_directional_event_study_high_discretion_income_added_back_500ft_paper.pdf}
    \caption{Permit Event Study: Unconstrained Directional Contrast}
    \figurenotes{The model estimates separate event-time coefficients for blocks assigned toward more stringent and more lenient aldermen.
    Each point is one-half of the difference between those coefficients.
    Shaded areas are 95\% confidence intervals based on standard errors clustered by ward pair.
    The pooled directional contrast is $-0.086$ (SE $=0.037$).
    The symmetry test for the pooled coefficients has $p=0.885$, and the joint pre-trend test has $p=0.657$.}
    \label{fig:permit_event_study_unconstrained}
\end{figure}

Figure~\ref{fig:permit_event_study_low_discretion} shows the full event-study path for low-discretion permits.
The response is more muted than for high-discretion permits.

\begin{figure}[htbp]
    \centering
    \includegraphics[width=0.7\textwidth]{../tasks/audits/permit_event_study_audit/output/binary_signed_directional_event_study_low_discretion_nosigns_income_added_back_500ft_paper.pdf}
    \caption{Event Study: Low-Discretion Permits}
    \figurenotes{The specification and sample match Figure~\ref{fig:permit_event_study}, but the outcome is low-discretion permits excluding signs, grouped by application year.
    Points report $\exp(\hat\beta_\tau)-1$; shaded areas are 95\% confidence intervals based on standard errors clustered by ward pair.
    The pooled estimate is $-0.043$ (SE $=0.033$; $p=0.200$).
    Pre-trend test $p=0.808$.}
    \label{fig:permit_event_study_low_discretion}
\end{figure}

\subsection{Stable-Incumbent Sample}
\label{sec:appendix_permit_stable}

As a robustness check, Figure~\ref{fig:permit_event_study_stable} retains only blocks whose 2014 origin- and destination-ward incumbents remained in office when the new map took effect.
The assigned and realized changes in stringency therefore coincide at implementation.
The estimated annual number of high-discretion permits is 11.6\% lower.
Because remaining in office is an election outcome, I treat this restriction as a robustness check rather than the main specification.

\begin{figure}[htbp]
    \centering
    \includegraphics[width=0.7\textwidth]{../tasks/audits/permit_event_study_audit/output/binary_signed_directional_event_study_high_discretion_income_added_back_500ft_stable_paper.pdf}
    \caption{Permit Event Study: Stable-Incumbent Sample}
    \figurenotes{The specification matches Figure~\ref{fig:permit_event_study}, but the sample retains only blocks whose 2014 origin- and destination-ward incumbents remained in office in 2015.
    Shaded areas are 95\% confidence intervals based on standard errors clustered by ward pair.
    The pooled estimate is $-0.124$ (SE $=0.046$), implying that the annual number of recorded permits is 11.6\% lower on average in years 0--5.
    Pre-trend test $p=0.386$.}
    \label{fig:permit_event_study_stable}
\end{figure}

\subsection{Permit Summary Statistics}
\label{app:permit_summary_stats}

Table~\ref{tab:permit_summary_stats} compares the high- and low-discretion permit groups defined in Section~\ref{subsec:permit_pipeline}.
They are much slower than lower-discretion permits: the mean processing time is 57.3 days versus 26.2 days, and the median is 30 days versus 6 days.
They also display more cross-alderman variation in processing times.

\begin{table}[htbp]
    \centering
    \caption{Summary Statistics for High- and Low-Discretion Permits}
    \label{tab:permit_summary_stats}
    \input{../tasks/permit_summary_stats/output/permit_processing_time_high_vs_low_alderman_summary.tex}

    \footnotesize
    \textit{Notes:} Sample uses positive-duration permit records with applications begun through December 2022.
    The high- and low-discretion permit groups are defined in Section~\ref{subsec:permit_pipeline}.
    The alderman-level interquartile range is computed from the cross-section of alderman-specific mean processing times within each permit group.
\end{table}

\subsection{Permit Volume and Processing Time}

Table~\ref{tab:permit_volume_correlation} reports the correlation between mean processing time and permit volume across ward-month cells.
The negative correlation for high-discretion permits is consistent with political and regulatory friction reducing throughput where discretionary review matters most, while the positive correlation for low-discretion permits suggests routine administrative scale rather than aldermanic discretion.

\begin{table}[htbp]
    \centering
    \caption{Processing Time and Permit Volume Correlation}
    \label{tab:permit_volume_correlation}
    \input{../tasks/permit_summary_stats/output/permit_processing_time_volume_correlation.tex}

    \footnotesize
    \textit{Notes:} Sample uses positive-duration permit records with applications begun through December 2022.
    Entries report the Pearson correlation between mean processing time and permit volume across ward-month cells for each permit group.
\end{table}
